Fix \(K\ge2\) and a contrast \(c\) satisfying \(\sum_ac_a=0\) and \(\sum_ac_a^2=1\). Define
\[
B_n=n,\qquad q_n=2n,\qquad n_{nb}=Kq_n=2Kn,
\qquad m_{nba}=q_n,\qquad w_{nb}=B_n^{-1}.
\tag{1}
\]
In every block label the units by \(i=1,\ldots,2Kn\), and set
\[
x_{nbi}=\begin{cases}-1,&i\le Kn,\\1,&i>Kn,\end{cases}
\qquad
Y_{nbi}(a)=x_{nbi},\qquad g_{nbi}(a)=0
\quad(a\in\mathcal A).
\tag{2}
\]
The subclass consists of this construction and arbitrary deterministic within-block permutations of the balanced list, using the same permutation in every arm of a block. In each block, conditional on the past, assign uniformly over allocations having exactly \(q_n\) units in every arm. Thus counts and adjustments are deterministic and predictable, all potential outcomes are fixed before randomization, consistency defines the observed response, and assignment is the only source of randomness.

We verify membership in one common-parameter class. Take
\[
\varepsilon_0=\frac1{2K},\qquad \delta_0=1,
\qquad C_R=1,\qquad v_0=K,\qquad \ell_n=\frac1n.
\]
Since \(p_{nba}=1/K\),
\[
\varepsilon_0\le p_{nba}\le1-(K-1)\varepsilon_0.
\]
The residuals equal \(x_{nbi}\), so
\(n_{nb}^{-1}\sum_i|R_{nbi}(a)|^{4+\delta_0}=1=C_R\). The minimum arm count is \(q_n=2n\to\infty\). Also \(N_n=B_nKq_n\), whence
\[
s_n^2=\sum_{b=1}^{B_n}\frac{w_{nb}^2}{n_{nb}}
=\frac1{B_nKq_n},
\qquad
\frac{w_{nb}^2/n_{nb}}{s_n^2}
=\frac1{B_n}=\ell_n\longrightarrow0.
\tag{3}
\]
The two contrast restrictions hold by construction. This verifies every member property of \(\mathrm{def:admissible\text{-}arrays}\) except predictable-QV growth, which is checked below.

The residual mean in every arm and block is zero, and the common finite-population variance is
\[
S_{nba}^2=S_n^2=\frac{Kq_n}{Kq_n-1}.
\tag{4}
\]
The diagonal coupling \((X_1,\ldots,X_K)=(X,\ldots,X)\), with \(X\) uniform on the balanced list in (2), lies in \(\Gamma_{nb}\). Since \(\sum_ac_a=0\), its cost is
\[
E\left(\sum_ac_aX_a\right)^2
=E\left(X\sum_ac_a\right)^2=0.
\]
The cost is nonnegative, so \(J_{nb}^{-}=0\). The true common-label coupling is this diagonal coupling. The conditional-variance identity in \(\mathrm{thm:exact\text{-}predictable\text{-}ot\text{-}envelope}\), together with \(\sum_ac_a^2=1\), gives
\[
E(D_{nb}^2\mid\mathcal F_{n,b-1})
=U_{nb}=\frac{S_n^2}{B_n^2q_n},
\qquad
V_n=U_n=\frac{S_n^2}{B_nq_n}.
\tag{5}
\]
This quantity is deterministic. Equations (3)--(5) imply
\[
\frac{V_n}{s_n^2}=KS_n^2\ge K=v_0,
\qquad
\frac{V_n}{E V_n}=1.
\tag{6}
\]
Thus predictable-QV growth and predictable-QV stabilization hold uniformly over the subclass, the latter exactly.

The smallest arithmetic witness is the case \(K=2\), \(n=1\), for which the block list is \(\{-1,-1,1,1\}\). With unnormalized contrast \((1,-1)\), diagonal pairing gives squared contrast zero at every atom, while reverse pairing gives squared contrast four at every atom. Thus the two transport costs are \(0\) and \(4\). With the required normalized contrast \((1,-1)/\sqrt2\), both costs are divided by two. This exactly checks the endpoint direction and confirms nontrivial width.

By \(\mathrm{thm:history\text{-}uniform\text{-}plugin}\),
\(\widehat U_n-U_n=o_{\Pr}(s_n^2)\) under the original randomization. Since (4)--(6) give \(U_n/s_n^2=KS_n^2\to K>0\),
\[
\frac{\widehat U_n}{V_n}=\frac{\widehat U_n}{U_n}\longrightarrow1
\quad\text{in probability}.
\tag{7}
\]

We verify the martingale CLT from the primitives. Estimator algebra, consistency, and \(p_{nba}=m_{nba}/n_{nb}\) give \(\widehat\tau_n-\tau_n=M_{nB_n}\). For an arm sample mean
\(X_a=q_n^{-1}\sum_{i:Z_{nbi}=a}x_{nbi}\), the fourth-moment expansion by inclusion multiplicities, with \(N=Kq_n\), \(A_2=A_4=N\), and \(\pi_j=(q_n)_j/(N)_j\), gives
\[
E(q_nX_a)^4
=\{\pi_1-7\pi_2+12\pi_3-6\pi_4\}A_4
+3\{\pi_2-2\pi_3+\pi_4\}A_2^2
\le26N+12N^2.
\]
Hence \(E(X_a^4\mid\mathcal F_{n,b-1})\le C_Kq_n^{-2}\). Since
\(D_{nb}=B_n^{-1}\sum_ac_aX_a\), \(\sum_ac_a^2=1\), and
\[
\left(\sum_ac_aX_a\right)^4\le K\sum_aX_a^4,
\]
we obtain
\[
\sum_{b=1}^{B_n}ED_{nb}^4
\le\frac{C_K}{B_n^3q_n^2},
\qquad
\frac{\sum_bED_{nb}^4}{V_n^2}
\le\frac{C_K}{B_n(S_n^2)^2}\longrightarrow0.
\tag{8}
\]
For every \(\epsilon>0\), the conditional Lindeberg sum divided by \(V_n\) has expectation at most \(\epsilon^{-2}\sum_bED_{nb}^4/V_n^2\), by \(u^2\mathbf1\{|u|>x\}\le u^4/x^2\). It converges to zero in probability by Markov's inequality. Together with \(V_n/EV_n=1\), the established martingale central limit theorem \(\mathrm{lem:martingale\text{-}clt\text{-}predictable\text{-}variation}\) gives
\[
T_n:=\frac{\widehat\tau_n-\tau_n}{\sqrt{V_n}}
\Rightarrow\mathcal N(0,1).
\tag{9}
\]
The causal estimand remains distinct from this randomization limit: indeed, \(\tau_n=0\), because every unit's fixed causal contrast is \(x_{nbi}\sum_ac_a=0\).

For fixed \(\rho\in(0,1)\), the shortened interval covers exactly when
\[
|T_n|\le z_{1-\alpha/2}
\sqrt{(1-\rho)\widehat U_n/V_n}.
\]
On the event \(|\widehat U_n/V_n-1|\le\zeta\), its coverage probability is sandwiched, up to the probability of the complementary event, between the probabilities that \(|T_n|\) is at most \(z_{1-\alpha/2}\sqrt{(1-\rho)(1-\zeta)}\) and \(z_{1-\alpha/2}\sqrt{(1-\rho)(1+\zeta)}\). Equations (7)--(9), continuity of \(\Phi\), and then \(\zeta\downarrow0\) yield
\[
\Pr\!\left\{|\widehat\tau_n-\tau_n|\le z_{1-\alpha/2}
\sqrt{(1-\rho)\widehat U_n}\right\}
\longrightarrow
2\Phi\!\left(z_{1-\alpha/2}\sqrt{1-\rho}\right)-1.
\tag{10}
\]
Because \(z_{1-\alpha/2}>0\), \(0<\sqrt{1-\rho}<1\), and \(\Phi\) is strictly increasing, the limit in (10) is strictly below \(2\Phi(z_{1-\alpha/2})-1=1-\alpha\).

Finally, if a broader design family contains this subclass, its worst-case coverage is no larger than coverage on the constructed member. Taking \(\limsup\) and applying (10) gives the stated matching fixed-fraction undercoverage certificate. The witness is nonadaptive and fixed-quota, so this inference necessity result does not assert cross-history compatibility of independently selected nodewise transport optimizers.